\section{Conclusion}

\model connects task-grounded security requirements to causal analysis and
evidence-grounded prompt guidance. Its structured representation locates the
requirement and the behavior to preserve; its causal analysis defines and
evaluates explicit prompt-policy comparisons; and its explanation contract
links the resulting evidence to a scoped modification decision. Observational
prioritization and independent randomized evaluation support this connection,
while retaining their distinct roles in acquiring effect evidence.

A policy-effect estimate can be informative without justifying adoption.
The explanation must retain the tested comparison, model and population,
measurement limits, functionality, and uncertainty, and withhold advice when
the required support is missing. Computational effects and expert-perceived
explanation utility address complementary questions; neither establishes the
other. The intended output is an inspectable connection between a proposed
requirement change, its supporting evidence, and its applicable scope, not a
universal security rule or an account of the generator's internal reasoning.
